\documentclass[12pt]{article}
\usepackage[a4paper,margin=1in]{geometry}
\usepackage{hyperref}
\usepackage{setspace}
\usepackage{titlesec}

\titleformat{\section}{\large\bfseries}{\thesection}{1em}{}

\onehalfspacing

\title{Project Report: Computer Science Portfolio}
\author{Efe Yalın Güleç}
\date{\today}

\begin{document}

\maketitle

\section{Introduction}

This report documents the design, implementation, and structure of my CV, created as part of the assessment task. The goal of the portfolio is to present myself as an applicant for a working student position and to showcase my technical skills, projects, and experience.

The portfolio is hosted using GitHub Pages and is publicly accessible at:

\begin{itemize}
    \item Portfolio Website: \url{https://theblue0n3.github.io/Efe-Yalin-Gulec-CV/}
    \item GitHub Repository: \url{https://github.com/TheBlue0n3/Efe-Yalin-Gulec-CV}
\end{itemize}

This report explains the structure of the website and repository, the design decisions behind the portfolio, the tools used to build it, and reflections on its strengths, weaknesses, and things i could do better later on.

\section{Portfolio Structure}

The portfolio consists of a single-page website built using HTML and CSS. The structure is intentionally minimal to ensure clarity, readability, and ease of navigation. The main sections of the website are:

\begin{itemize}
    \item \textbf{About Me} – A short introduction describing my background and interests.
    \item \textbf{Skills} – A summary of the programming languages and tools I am familiar with.
    \item \textbf{Projects} – Links to selected GitHub repositories showcasing my work.
    \item \textbf{CV} – Pretty Self explanatory.
    \item \textbf{Contact} – Links and information for reaching me.
\end{itemize}

The website uses a dark-mode, for the better of everybody's eye health. Make sure you use yellow light when your enviorment is dark everyone.

\section{Repository Structure}

The GitHub repository contains all files required for the portfolio website and supporting materials. The structure is as follows:

\begin{itemize}
    \item \texttt{index.html} – The main webpage.
    \item \texttt{README.md} – Documentation describing the repository.
    \item \texttt{Efe\_Yalin\_Gulec\_Resume.pdf} – My LaTeX-generated CV and everything else related to it.
\end{itemize}

Larger projects are not duplicated inside this repository; instead, they are linked directly to their respective GitHub repositories. Neat.

\section{Design Decisions}

Several design choices were made to ensure the portfolio is both functional and visually appealing:

\begin{itemize}
    \item \textbf{Dark Mode:} Chosen for readability, modern aesthetics, and reduced eye strain.
    \item \textbf{Minimal Layout:} A simple structure helps hiring managers quickly find relevant information.
    \item \textbf{HTML + CSS Only:} Keeping the site lightweight ensures fast loading and easy maintenance.
    \item \textbf{GitHub Pages Hosting:} Provides free, reliable hosting directly integrated with version control.
    \item \textbf{Project Linking:} Instead of embedding full projects, links to GitHub repositories maintain clarity and avoid clutter.
\end{itemize}

Everything that a professional and nice looking CV needs.

\section{Tools and Technologies Used}

The following tools and technologies were used to develop the portfolio:

\begin{itemize}
    \item \textbf{HTML/CSS} – For building and styling the webpage.
    \item \textbf{Git and GitHub} – For version control, hosting, and project management.
    \item \textbf{GitHub Pages} – For deploying the portfolio website.
    \item \textbf{LaTeX} – For creating both this report and my CV.
    \item \textbf{VS Code / Overleaf} – For editing code and LaTeX documents.
\end{itemize}

These tools were selected because they are widely used in industry and what my colleagues largely recommended to me.

\section{Reflection}

\subsection*{Strengths}

The portfolio has several strengths:

\begin{itemize}
    \item Clean and modern design.
    \item Easy-to-navigate structure.
    \item Clear presentation of skills and projects.
    \item Professional use of GitHub Pages and version control.
    \item Lightweight and fast-loading.
    \item Allows me to more easily get employed
\end{itemize}

\subsection*{Weaknesses}

There are also areas that could be improved:

\begin{itemize}
    \item Limited visual elements and interactivity.
    \item No responsive design for mobile devices.
    \item No navigation bar for quick section access.
    \item Can't bake you a cake, I tried.
\end{itemize}

\subsection*{Future Improvements}

To further enhance the portfolio, I plan to:

\begin{itemize}
    \item Add responsive design for better mobile usability.
    \item Implement a navigation bar for smoother navigation.
    \item Include more detailed project descriptions and visuals.
\end{itemize}

\section{Conclusion}

This portfolio demonstrates my technical skills, projects, and ability to design and deploy a functional website using industry-standard tools. The process of building the portfolio and writing this report has strengthened my understanding of web development, and documentation. 

\end{document}
